\begin{problem}[Warmup]{(a)}
    Compute explicitly three conjugates of the permutation $(145)(96)(237)$, checking that its disjoint cycle structure is preserved.
\end{problem}
\unsolved

\begin{problem}[Warmup]{(b)}
    How many elements are in $A_5$? $A_6$?
\end{problem}
\unsolved

\begin{problem}[Warmup]{(c)}
    Because conjugation preserves parity, the conjugate of an odd permutation is odd. Why doesn't that make the set of odd permutations a normal subgroup?
\end{problem}
\unsolved

\begin{problem}[Warmup]{(d)}
    Why can't a simple group have a non-trivial subgroup of index 2?
\end{problem}
\unsolved

\begin{problem}[Warmup]{(e)}
    Give an example of a non-normal subgroup, whose index is 3. (Thus, Theorem 32.3 is false if 2 is replaced by 3.)
\end{problem}
\unsolved

\begin{problem}[Warmup]{(f)}
    Is the identity permutation even or odd?
\end{problem}
\unsolved

\begin{problem}[Warmup]{(g)}
    Suppose the operation $\beta$ means ``paint the north wall of your room'', and the operation $\alpha$ means ``move to your sister's room''. What is the meaning of the conjugate $\alpha^{-1}\beta\alpha$? Remember that functional composition is read from right to left.
\end{problem}
\unsolved

\begin{problem}{1}
    Is the product of an even permutation and an odd permutation always an odd
    permutation? Prove this, or give a counterexample.
\end{problem}
\unsolved

\begin{problem}{2}
    Show that the map $\Lambda : A_n \to S_n \setminus A_n$ defined by
    $\Lambda(\alpha) = (12)\alpha$ is one-to-one and onto, completing the argument of
    Section 34.3 that there are as many even as odd permutations. (The notation
    $S_n \setminus A_n$ means the elements of $S_n$ that are not in $A_n$.)
\end{problem}
\unsolved

\begin{problem}{3}
    \begin{enumerate}
        \item[(a)] List the elements of $A_5$. (See Exercise b.)
        \item[(b)] List all cyclic subgroups of $A_5$.
        \item[(c)] Let $\alpha \in A_5$. Show $\langle \alpha \rangle$ is not a normal
        subgroup of $A_5$.
    \end{enumerate}
\end{problem}
\unsolved

\begin{problem}{4}
    Show that $A_n$ is simple for $n < 4$, proving part of Theorem 34.4.
\end{problem}
\unsolved

\begin{problem}{5}
    By following the proof of the Conjugation Theorem 34.2, demonstrate explicitly that
    \[(1456)(29) \quad\text{and}\quad (7895)(13)\]
    are conjugates in $S_9$.
\end{problem}
\unsolved

\begin{problem}{6}
    Let $n \ge 5$. Prove that $A_n$ is the only non-trivial normal subgroup of $S_n$.
\end{problem}
\unsolved

\begin{problem}{7}
    Prove that any permutation can be expressed as a product of the transpositions
    $(12), (13), (14), \ldots, (1n)$.
\end{problem}
\unsolved

\begin{problem}{8}
    Prove that any permutation can be expressed as a product of the transpositions
    $(12), (23), (34), \ldots, (n - 1\ n)$.
\end{problem}
\unsolved

\begin{problem}{9}
    Prove that $S_n$ is generated by the two elements $(1234\cdots n)$ and $(12)$.
\end{problem}
\unsolved

\begin{problem}{10}
    In Example 31.6 we showed that $A_4$ has no subgroup with 6 elements. Provide a
    different proof for this fact: Assume that $H$ is such a subgroup; then $H$ must be
    normal. (Why?) Now count the three cycles. Show that at least one of the three cycles
    must be in $H$. Then argue that $H$ must contain all the three cycles, which is absurd.
\end{problem}
\unsolved

\begin{problem}{11}
    Let $n \ge 2$ and let $H$ be a subgroup of $S_n$. Prove either that all elements of
    $H$ are even, or else exactly half of the elements of $H$ are even.
\end{problem}
\unsolved

\begin{problem}{12}
    Let $n$ be a positive integer, $n \ge 3$. Prove that $A_n$ is generated by its
    3-cycles. That is, prove that every element of $A_n$ can be expressed as a product of
    3-cycles.
\end{problem}
\unsolved

\begin{problem}{13}
    Suppose that $H$ is a subgroup of the permutation group $S_n$. We say that $H$ is a
    \textbf{transitive subgroup} if for any integers $k, j$ with $1 \le k, j \le n$, we can
    find $\alpha \in H$ so that $\alpha(k) = j$.
    \begin{enumerate}
        \item[(a)] Determine all transitive subgroups of $S_4$.
        \item[(b)] Argue that $A_n$ is a transitive subgroup of $S_n$, for any $n > 2$.
        \item[(c)] Suppose that $H$ is a transitive subgroup of $S_n$, and
        $\alpha \in S_n$. Prove that the conjugate subgroup $\alpha^{-1}H\alpha$ is also a
        transitive subgroup of $S_n$. (See Exercise 32.6 for more about conjugate
        subgroups.)
    \end{enumerate}
\end{problem}
\unsolved
